\section{Requirements for a New Dataset}
\label{sec:dataset-requirements}

Answering the research questions defined in Section~\ref{sec:rq-derivation} requires both a single-turn and a multi-turn dataset. In this work, a \emph{sample} is a single data point of a dataset: one request, together with its target label. In both datasets, each sample must be a request meant to be submitted to a coordinator agent, which must route it to one of its available sub-agents. Each sample must also carry a target label, which defines the sub-agent to which the request must be routed. The rest of this section describes each requirement.

\subsection{Routing Task}
\label{subsec:routing-setup}

The number of sub-agents used must satisfy the following requirements.

\begin{itemize}
    \item It must stay within the range Section~\ref{sec:multi-agent} reports for a production coordinator, roughly three to seven sub-agents.
    \item It must allow requests for different sub-agents to share vocabulary, so that the wording alone does not reveal which sub-agent must handle the request.
\end{itemize}

Each sample must still be routable to one, and only one, sub-agent, so that a prediction on that sample can be marked right or wrong, depending on whether the chosen sub-agent matches the correct one. The shared vocabulary must therefore hold at the level of the wording, while the content of the request still selects a single sub-agent.

\subsection{Difficulty for Text-Only Classifiers}
\label{subsec:lexical-separability}

To answer \hyperref[rq:rq2]{\textbf{RQ2}}, the dataset must be challenging for text-only classifiers. This gives the probe margin to score above them.

Three text-only classifiers are used to measure how well this difficulty holds: a logistic regression trained on bag-of-words counts, a logistic regression trained on TF-IDF-weighted counts, and a logistic regression trained on sentence embeddings. The dataset must aim for the lowest accuracy these classifiers can reach, while every sample still reads as natural and remains routable to exactly one sub-agent, without an ambiguity that would leave it unroutable.

The sub-agents must also have a balanced number of requests, and be balanced in writing style, punctuation, and sentence pattern, in addition to vocabulary. Otherwise, a text-only classifier could route a request by its class frequency or by a stylistic marker instead of by what the request asks.

\subsection{Multi-turn Setting}
\label{subsec:multiturn-setting}

The multi-turn dataset must be built on the same sub-agents as the single-turn dataset, with the same writing styles, the same number of samples, and the same balance of samples across sub-agents. The only difference is that each multi-turn sample distributes its request across more than one turn, as defined in Chapter~\ref{chap:llm} (Section~\ref{subsec:query-prompt}), instead of stating it in one.

Each multi-turn sample must use the same number of turns, so that the length of the conversation is controlled. This count must stay fixed because the dataset is meant to show how performance changes between the single-turn and the multi-turn setting, not how it changes with the number of turns.

Each multi-turn sample must start with a request and end with a request. The final request must have no response: it is the request to be routed. Together, the turns in a sample are required to carry enough information to route the complete conversation to exactly one sub-agent.
